% Sections 1.1 to 1.6, from lectures/L01/appendix/a_combinational_logic.md: one section per part,
% A.1 to A.6.

\section{Digital Signals and Logic Gates}
\label{c1:a1}
A digital signal takes one of two values, \code{0} or \code{1}, rather than the continuous range an
analog signal uses. Anything close enough to \code{0} or \code{1} is read as exactly \code{0} or
\code{1}, so small amounts of electrical noise change nothing. That noise immunity is why almost all
modern computation, from microcontrollers to FPGAs, is digital.

A \textbf{logic gate} has one or more single-bit inputs, one single-bit output, and a fixed Boolean
function between them.

\textbf{Combinational logic}, this chapter's topic, produces outputs from the current inputs alone,
with no memory. \textbf{Sequential logic} (Chapter~\ref{c3:ch}) depends on the current inputs
\emph{and} on state stored in flip-flops.

Every basic 2-input gate, plus the unary \code{NOT}:

\begin{narrowtable}{@{}cccccccc@{}}
  \thead{A} & \thead{B} & \thead{AND} & \thead{OR} & \thead{NAND} & \thead{NOR} & \thead{XOR} &
    \thead{XNOR} \\
  \midrule
  0 & 0 & 0 & 0 & 1 & 1 & 0 & 1 \\
  0 & 1 & 0 & 1 & 1 & 0 & 1 & 0 \\
  1 & 0 & 0 & 1 & 1 & 0 & 1 & 0 \\
  1 & 1 & 1 & 1 & 0 & 0 & 0 & 1 \\
\end{narrowtable}

\begin{narrowtable}{@{}cc@{}}
  \thead{A} & \thead{NOT A} \\
  \midrule
  0 & 1 \\
  1 & 0 \\
\end{narrowtable}

\noindent Worth internalizing straight from the tables:
\begin{itemize}
  \item \code{NAND}, \code{NOR} and \code{XNOR} are the inverses of \code{AND}, \code{OR} and
    \code{XOR}.
  \item \code{XOR} is \code{1} when the inputs differ. Beyond two inputs this generalizes to
    \textbf{odd parity}: \code{1} when an odd number of inputs are \code{1}
    (Exercise~\ref{c1:ex:3}).
  \item Any Boolean function can be built from \code{AND}, \code{OR} and \code{NOT}, and also from
    \code{NAND} alone or \code{NOR} alone.
\end{itemize}

\section{Boolean Algebra Notation}
\label{c1:a2}
Three operators, mapping directly onto the gates above:

\begin{narrowtable}{@{}ll@{}}
  \thead{Gate} & \thead{Boolean algebra notation} \\
  \midrule
  AND & multiplication: \code{A} \textbf{AND} \code{B} is written \code{AB} \\
  OR & addition: \code{A} \textbf{OR} \code{B} is written \code{A + B} \\
  NOT & prime: \textbf{NOT} \code{A} is written \code{A'} \\
\end{narrowtable}

\noindent A function written as an OR of AND-terms, such as \code{X = AB' + CD}, is in
\textbf{sum-of-products} form.

Every truth table reads off into sum-of-products directly:
\begin{enumerate}
  \item Find every row where the output is \code{1}.
  \item For each, write one AND-term containing every input: directly if it is \code{1} in that
    row, primed if it is \code{0}.
  \item OR those terms together.
\end{enumerate}
For example:

\begin{narrowtable}{@{}ccc@{}}
  \thead{A} & \thead{B} & \thead{X} \\
  \midrule
  0 & 0 & 0 \\
  0 & 1 & 1 \\
  1 & 0 & 1 \\
  1 & 1 & 0 \\
\end{narrowtable}

\noindent Row \code{AB = 01} contributes \code{A'B} and row \code{AB = 10} contributes \code{AB'},
giving \code{X = A'B + AB'}. Check it against the table: that is \code{XOR}, so the same network is
one \code{XOR} gate rather than two \code{AND} gates, two \code{NOT} gates and an \code{OR}.

\subsection{The laws, as a reminder}
\label{c1:laws}
You have met these. They are here so the book has one notation for them and so the later chapters
can cite them, not because they need teaching:

\begin{narrowtable}{@{}lll@{}}
  \thead{Law} & \thead{With OR} & \thead{With AND} \\
  \midrule
  Identity & \code{A + 0 = A} & \code{A · 1 = A} \\
  Null & \code{A + 1 = 1} & \code{A · 0 = 0} \\
  Idempotent & \code{A + A = A} & \code{A · A = A} \\
  Complement & \code{A + A' = 1} & \code{A · A' = 0} \\
  Absorption & \code{A + AB = A} & \code{A(A + B) = A} \\
  De Morgan & \code{(A + B)' = A'B'} & \code{(AB)' = A' + B'} \\
\end{narrowtable}

\noindent\textbf{De Morgan is the one this book leans on}, less for minimization than as a rule
about moving an inversion through a gate: an \code{OR} with both inputs inverted \emph{is} a
\code{NAND}, and an \code{AND} with both inputs inverted \emph{is} a \code{NOR}. It is why a
\code{NAND} gate can build anything, which Exercise~\ref{c1:ex:6} has you write in VHDL, and it
comes back every time an active-low signal meets logic written for active-high. In a book where
every reset and every button is active-low, that is often.

That gap between a straightforward sum-of-products reading and the smallest network is what a
\textbf{Karnaugh map} closes, and Chapter~\ref{c2:ch} works one end to end. Not because the
technique is new to you, but because the minimized equation is what gets written in VHDL, and
Chapter~\ref{c8:ch} derives its state machine's next-state logic exactly this way.

\section{Building and Simulating a Circuit in CircuitVerse}
\label{c1:a3}
CircuitVerse (\url{https://circuitverse.org/simulator}) is a free, browser-based logic simulator.
Every gate network in this book is built and tested there by hand before it is written in VHDL:
\begin{enumerate}
  \item Open the simulator and start a new project.
  \item Place an \textbf{input} element per input signal and an \textbf{output} element per output,
    renamed to match your variables.
  \item Place the gates your equation calls for, from the \emph{Gates} panel.
  \item Wire it up one gate at a time, following your derived equation. You should already know
    exactly which gates you need before you open the simulator, which is why \secref{c1:a2} comes
    first.
  \item Toggle each input and confirm the output matches your truth table for \textbf{every} row.
  \item Save via \code{Project -> Save Online}, or \code{Project -> Download as image/data} for a
    re-importable \file{.cv} file.
\end{enumerate}
\textbf{Step 5 is not optional, and nothing else does it for you.} Three rules make it a check
rather than a formality:
\begin{itemize}
  \item \textbf{Exhaustive, not spot-checked.} Two inputs is four rows, three is eight. At these
    sizes there is no excuse for sampling. For the sequential circuits later on this becomes ``step
    the clock and check at every edge'', the same discipline applied to time.
  \item \textbf{Derive first, simulate second.} Draw the circuit and \emph{then} decide what it
    should do, and you will read your expectations off the simulation, which proves nothing.
  \item \textbf{A mismatch is information.} It says the drawing and the equation disagree, and one
    of them is wrong.
\end{itemize}
Every hardware bug you will chase is found by predicting a behaviour and then checking it. Doing
it by hand on a four-gate circuit is how you learn to do it on one you cannot see all at once.

Once the design is VHDL, a self-checking testbench takes over that job. Every worked example and
every exercise ships one; you run them rather than write them, as covered in \secref{c2:app:c}.

\section{Enough VHDL to Read and Write a Gate Circuit}
\label{c1:a4}
VHDL describes hardware, not a sequence of instructions. Every statement below runs concurrently
and continuously, exactly like the gates it describes. The language arrives as it is needed; this
section is enough for a simple combinational circuit.

Every VHDL module has two parts:
\begin{itemize}
  \item An \textbf{entity}, the black-box view: the ports, and nothing about what happens inside.
  \item An \textbf{architecture}, the implementation behind it.
\end{itemize}
That split is a contract, not ceremony. The entity is everything another design needs in order to
use this module, which is what lets one design instantiate another without reading it, and lets
you rewrite an implementation without touching what depends on it. The shape is a header beside a
source file, or an interface beside its class; the difference is that VHDL enforces it for every
module, with no way to skip the interface or leak the implementation through it.

Take \code{or_gate}: inputs \code{a}, \code{b}, output \code{x}, and \code{x = a or b}.

\begin{vhdlcode}
entity or_gate is
    port(a, b: in std_logic;
         x   : out std_logic);
end entity;
\end{vhdlcode}

\bookfigure{or_gate_entity}{The entity \code{or_gate}: two inputs and an output, and nothing about
what happens inside.}{c1:fig:entity}

\noindent\code{std_logic} is the single-bit signal type. It is not built into VHDL but comes from
the \code{std_logic_1164} package, which is why every file in this book opens with:

\begin{vhdlcode}
library ieee;
use ieee.std_logic_1164.all;
\end{vhdlcode}

\noindent It has nine values rather than two, because real hardware needs more than true and false:
\begin{itemize}
  \item \code{'0'} and \code{'1'}: driven low and driven high. The only two you will write in this
    book.
  \item \code{'Z'}: high-impedance, nothing driving the line at all.
  \item \code{'U'}: uninitialized, what a signal holds before anything assigns it.
  \item \code{'X'}: unknown, what a simulator shows when two things drive one signal and disagree.
  \item \code{'-'}, \code{'W'}, \code{'L'}, \code{'H'}: don't-care and weak-drive values, rarely
    written by hand.
\end{itemize}
Recognizing the seven you will not write still matters: they are how a simulator tells you
something is wrong. A signal sitting at \code{'U'} or \code{'X'} is a bug, not a value.

The architecture fills in the box:

\begin{vhdlcode}
architecture behaviour of or_gate is
begin
    x <= a or b;
end architecture;
\end{vhdlcode}

\bookfigure{or_gate_arch}{The architecture of \code{or_gate}: the gate that fills the
box.}{c1:fig:arch}

\noindent Together they form the complete module:

\bookfigure{or_gate_module}{The module \code{or_gate}: the entity and the architecture
together.}{c1:fig:module}

\noindent\code{x <= a or b;} is \secref{c1:a2}'s Boolean algebra almost verbatim: VHDL's
\code{and}, \code{or}, \code{not} and \code{xor} work directly on \code{std_logic}, so any equation
you derived by hand becomes VHDL with barely any translation.

The assignment is \textbf{concurrent}. It is not executed once; it holds continuously, like a wire
between real gates.

The keyword \code{signal} declares an internal wire inside an architecture, as opposed to an
externally visible port. You use one in Exercise~\ref{c1:ex:7} to name an intermediate result
rather than nesting it inline, and Chapter~\ref{c2:ch} leans on the same idea throughout.

\section{The Complete Module: \code{or_gate}}
\label{c1:a5}
\Secref{c1:a4}'s pieces together make a complete, synthesizable file. This is the whole of
\file{lectures/L01/or_gate/or_gate.vhd}, and it is a real design that synthesizes and runs on the
DE0-CV:

\vhdlfile{lectures/L01/or_gate/or_gate.vhd}

\begin{itemize}
  \item The \code{library}/\code{use} clauses pull in \code{std_logic}.
  \item \code{entity} declares the outside: two inputs, one output.
  \item \code{architecture} implements the inside, as one concurrent assignment.
\end{itemize}
That is the smallest complete VHDL design there is, and the shape never changes. Every module in
this book, up to the state machines in Chapter~\ref{c8:ch}, is this same
\code{library}/\code{entity}/\code{architecture} skeleton. What grows is the architecture body,
never the frame.

\textbf{Checking your work.} \code{or_gate} ships \file{or_gate_tb.vhd}, which drives all four
input combinations and checks the output:

\begin{shell}
cd lectures/L01/or_gate
ghdl -a --std=93 or_gate.vhd or_gate_tb.vhd
ghdl -e --std=93 or_gate_tb
ghdl -r --std=93 or_gate_tb --assert-level=error --stop-time=10ms
\end{shell}

\noindent No output beyond a final note means every check passed. Testbenches are covered properly
in \secref{c2:app:c}; for now, treat these three commands as how you confirm a module works.

\section{The Worked Circuit: A Braking Assistant}
\label{c1:a6}
The course builds a different circuit from \code{or_gate}, live, so you see the same path walked
twice: once here in writing, and once on something you have not already read the answer to.

The requirement: \textbf{the car brakes if the driver brakes, or if the assistance system decides
to.} The assistance system decides to when either detector sees something, unless it is reporting a
fault, in which case it is not trusted to brake at all.

\begin{booktable}{@{}p{7.5em}p{5em}L@{}}
  \thead{Port} & \thead{Direction} & \thead{Meaning} \\
  \midrule
  \code{driver_brake} & in & The driver is on the brake pedal. \\
  \code{sensor} & in & The proximity sensor sees an obstacle. \\
  \code{radar} & in & The radar sees an obstacle. \\
  \code{error} & in & The assistance system is reporting a fault. \\
  \code{engine_brake} & out & Apply the brakes. \\
\end{booktable}

\noindent Three clauses decide the whole structure between them:
\begin{itemize}
  \item either detector is enough on its own, and neither is trusted more than the other.
  \item a fault does not raise an alarm; it \emph{suppresses} the assistance system's ability to
    brake.
  \item the driver can always brake, including when the assistance system is faulty and doing
    nothing.
\end{itemize}
That last one is a safety requirement, and it is the one to hold on to when you derive the
network: the pedal is not an input to the assistance logic, it bypasses it. Work out what changes
if the pedal is fed through the fault logic instead, and you have found the bug the structure
exists to prevent.

\subsection{The truth table}
\label{c1:adastable}
Sixteen rows, in the order a four-bit counter would produce them. \code{adas_brake} is not a port;
it is the assistance system's own decision, shown as an intermediate column so you can see where
the structure comes from.

\begin{narrowtable}{@{}cccccc@{}}
  \thead{\code{driver_brake}} & \thead{\code{sensor}} & \thead{\code{radar}} & \thead{\code{error}}
    & \thead{\code{adas_brake}} & \thead{\code{engine_brake}} \\
  \midrule
  0 & 0 & 0 & 0 & 0 & 0 \\
  0 & 0 & 0 & 1 & 0 & 0 \\
  0 & 0 & 1 & 0 & 1 & 1 \\
  0 & 0 & 1 & 1 & 0 & 0 \\
  0 & 1 & 0 & 0 & 1 & 1 \\
  0 & 1 & 0 & 1 & 0 & 0 \\
  0 & 1 & 1 & 0 & 1 & 1 \\
  0 & 1 & 1 & 1 & 0 & 0 \\
  1 & 0 & 0 & 0 & 0 & 1 \\
  1 & 0 & 0 & 1 & 0 & 1 \\
  1 & 0 & 1 & 0 & 1 & 1 \\
  1 & 0 & 1 & 1 & 0 & 1 \\
  1 & 1 & 0 & 0 & 1 & 1 \\
  1 & 1 & 0 & 1 & 0 & 1 \\
  1 & 1 & 1 & 0 & 1 & 1 \\
  1 & 1 & 1 & 1 & 0 & 1 \\
\end{narrowtable}

\noindent The bottom half is entirely \code{1}: once \code{driver_brake} is set nothing else changes
the answer. That is what ``overrides everything'' looks like in a table.

\subsection{Before you read on}
\label{c1:adasbefore}
Everything above is the specification. The derivation, the gate network, and the VHDL are what the
course builds live, and it is worth arriving at them with your own answer to compare against rather
than one you have already read. So, from the table above:
\begin{itemize}
  \item Read a sum-of-products equation for \code{engine_brake} straight off it, the way
    \secref{c1:a2} did.
  \item Simplify it by hand. The requirement is three clauses long, so the minimal network is a
    great deal smaller than the row-by-row reading gives you. Work out how much smaller before you
    are told.
  \item Count the gates your simplified equation needs, and write the number down.
  \item Sketch the network, and decide which single gate the safety argument above rests on.
\end{itemize}
Write those four answers down. Being wrong about any of them is worth more than not having guessed,
because you will know exactly which step to re-examine when you check them. That is
\secref{c1:a3}'s ``derive first, simulate second'' applied to a circuit you have not built yet, and
it is the habit the rest of the book is built on. Afterwards you write the VHDL yourself, in
Exercise~\ref{c1:ex:7}, and a testbench decides whether your equations and this truth table
agree.

\subsection{Why there is a testbench for it}
\label{c1:adastb}
Sixteen combinations is few enough to check by hand once. It is not few enough to check every time
somebody edits the file, which is the situation every module in this book is in from here on.

A \textbf{testbench} is a second VHDL file that does the checking. It is not part of the design and
never reaches the FPGA: it drives the inputs through every case, works out what the output should
have been, compares, and reports. For \code{adas}, the braking assistant's module, that is a loop
over the sixteen combinations restating the braking rule independently.

\emph{Independently} is the important word. A testbench that read its expected answer out of the
design would agree with it by construction, including when the design is wrong. Restating the
requirement separately is what turns it into a check.

That is why almost every directory in the companion repository has a \file{_tb.vhd} beside the
module: it is how you can be handed a module, change it, and know within a second whether you broke
it. \Secref{c2:app:c} covers running them; you are not asked to write one anywhere in this book.
